% arXiv-ready draft. Compile: pdflatex main && bibtex none (thebibliography inline) && pdflatex main
\documentclass[11pt]{article}
\usepackage[margin=1in]{geometry}
\usepackage{booktabs}
\usepackage{multirow}
\usepackage{url}
\usepackage{hyperref}
\usepackage{microtype}

\title{DiskANN-RSI: Query-Path Optimizations for a Rust-Native Vamana Index,\\and a Same-Machine Reproducibility Study}

\author{
  Zhanhui Kang\\
  Independent Researcher\\
  \texttt{kangzhanh@gmail.com} \quad \url{https://github.com/kangzhanhui}
}

\date{September 2026}

\begin{document}
\maketitle

\begin{abstract}
We present \textbf{DiskANN-RSI}, a Rust-native implementation of the Vamana/DiskANN
graph index whose query path was iteratively optimized by an automated
research-loop (RSI) framework guided by recent graph-ANN theory
(Gollapudi et al., ICML 2025). Under the optimizer's own evaluation harness,
the accumulated query-path changes raise single-thread
QPS-at-recall\,$\ge 0.95$ on SIFT1M from 2994.6 to 4641.1 (+55\%), with the
dominant single change being a cache-resident epoch-stamped visited set using
16-bit stamps. We then re-evaluate the same binary against unmodified
baselines (stock Vamana/DiskANN, hnswlib in three configurations, and NSG)
under one protocol on one machine. The isolated gains do not transfer
uniformly: DiskANN-RSI trails stock Vamana by 9--31\% at low search depths and
reaches parity at high recall, while building its index $1.9\times$ faster
(131.9s vs.\ 255.0s) and scaling slightly better to two threads (1.87$\times$
vs.\ 1.75$\times$). We analyze why the two evaluation protocols disagree,
report several negative results (SQ8 quantization, candidate-queue merge
insertion, random-projection reordering, partial-distance early exit), and
release code, raw measurements, and runbooks to support reproduction. Our
findings quantify a familiar but under-documented hazard: optimizer-in-the-loop
harnesses reward changes that do not survive protocol unification, and
same-machine cross-system baselines remain necessary.
\end{abstract}

\section{Introduction}
Approximate nearest neighbor (ANN) search on graph indexes underlies modern
vector retrieval. The Vamana graph, introduced with DiskANN
\cite{subramanya2019diskann}, remains a reference design for single-node
indexes, and recent work continues to tighten its construction guarantees
\cite{gollapudi2025sort}. Separately, automated research-loop systems
(``AI scientists'') \cite{lu2024aiscientist,yamada2025aiscientistv2} have
begun to propose, implement, and measure code optimizations with minimal
human involvement, raising a practical question: \emph{do gains reported by an
optimizer's own harness survive contact with a unified, same-machine
benchmark against independent baselines?}

This paper is both a system description and a reproducibility study.
We ported DiskANN's full component set to Rust (graph index, SSD-resident
index with aligned reads, PQ/SQ quantization, label filtering, streaming
executor), and let an automated research-and-iteration loop (RSI) optimize
the in-memory Vamana query path on SIFT1M. We then froze the resulting binary
and benchmarked it against stock Vamana, hnswlib, and NSG on one cloud host
under one protocol. Our contributions:

\begin{itemize}
  \item A Rust-native Vamana/DiskANN implementation with a reproducible,
  configuration-driven benchmark harness reporting QPS-at-recall operating
  points (Section~\ref{sec:system}).
  \item A set of query-path optimizations --- an epoch-stamped visited set
  with 16-bit stamps, adaptive early termination, exposed tuning levers, and
  fat link-time optimization (LTO) --- with per-change measurements
  (Section~\ref{sec:opts}).
  \item A same-machine, single-protocol comparison against stock Vamana,
  hnswlib (three variants), and NSG, including recall--QPS ladders, thread
  scaling, build time, and peak memory (Section~\ref{sec:eval}).
  \item An analysis of where the two evaluations diverge and a documented set
  of negative results (Sections~\ref{sec:analysis}--\ref{sec:negative}),
  with all raw data released.
\end{itemize}

\section{Background}
\label{sec:background}
\paragraph{Graph ANN.} HNSW \cite{malkov2020hnsw} builds a hierarchical
navigable small-world graph; NSG \cite{fu2019nsg} sparsifies a kNN graph
toward the monotonic-relative-neighborhood property; DiskANN's Vamana
\cite{subramanya2019diskann} uses greedy search plus a robust pruning rule
with a slack parameter $\alpha$, supporting both in-memory and SSD-resident
operation. Worst-case analyses of popular implementations
\cite{indyk2023worstcase} and of the DiskANN family
\cite{gollapudi2025sort} motivate construction-time changes
(e.g., sorted $\alpha$-reachable graphs); our work instead targets the
\emph{search-time} hot path.

\paragraph{Query-path costs.} A Vamana query is a best-first beam search:
pop the closest frontier vertex, compute distances to its out-neighbors,
mark them visited, and merge qualifying candidates into the frontier.
Costs concentrate in (i) distance computation, (ii) visited-set membership,
and (iii) frontier maintenance. Each is a known lever: visited sets are
classically implemented as hash tables or epoch-stamped dense arrays;
early termination heuristics trade recall for fewer distance computations;
quantized traversal with full-precision reranking reduces bandwidth
pressure when the working set exceeds cache.

\section{System Overview}
\label{sec:system}
The implementation is a Cargo workspace of 15+ crates in four tiers:
SIMD/vector primitives (\texttt{diskann-wide}, \texttt{diskann-vector}),
core math and utilities (\texttt{diskann-linalg}, \texttt{diskann-utils},
\texttt{diskann-quantization}), the algorithm and storage layer
(\texttt{diskann}, \texttt{diskann-providers}, \texttt{diskann-disk},
\texttt{diskann-label-filter}), and benchmarking/tooling
(\texttt{diskann-benchmark*}, \texttt{diskann-tools}). The benchmark runner
is driven by a declarative \texttt{job.json} and supports a
\texttt{recall\_target} summary: the best QPS whose recall meets the target,
reported per thread count --- the industry-standard QPS-at-recall operating
point used throughout this paper. Data preparation for SIFT1M is scripted
(\texttt{tools/prepare\_sift1m.py}); all measurements in
Section~\ref{sec:eval} use recall@10 over 10{,}000 queries with 3
repetitions, and we report medians.

\section{Query-Path Optimizations}
\label{sec:opts}
The optimizations below were proposed, implemented, and measured by the RSI
loop between 2026-09-14 and 2026-09-15. Unless noted, percentages are
measured under the RSI harness on SIFT1M (1M $\times$ 128d) with
QPS-at-recall\,$\ge 0.95$ as the primary metric.

\paragraph{O1: Epoch-stamped visited set, 16-bit stamps.}
We replace the hash-set visited structure with a dense array stamped per
query; clearing is an $O(1)$ epoch bump, with an overflow set for sentinel
ids. The first version (32-bit stamps) gave +10--15\% single-thread QPS.
Narrowing stamps to 16 bits halves the footprint (4\,MB $\to$ 2\,MB at 1M
points) for cache residency; instrumented breakdowns show visited-set time
$-45\%$ and beam-expansion time $-19\%$. Cumulative effect of the query-path
package: 1T QPS@0.95 $2994.6 \to 4641.1$ (+55\%), with latency down across
the ladder (L$=35$: $\sim$330\,$\mu$s $\to$ 215\,$\mu$s; L$=100$:
$\sim$480\,$\mu$s $\to$ 318\,$\mu$s) at identical recall and distance-computation
counts; two-thread scaling improves to 1.83$\times$ (8486.6 QPS@0.95).

\paragraph{O2: Adaptive early termination.}
An optional stop rule halts expansion when the frontier exceeds the current
$k$-th distance by a configurable ratio after a minimum hop count. On SIFT1M
at ratio 1.2 (committed default) it cuts $\sim$10\% of distance computations
at L$=60$ and $\sim$50\% at L$=200$ with recall statistically unchanged.
A finer sweep showed ratio 1.15 dominates near the recall knee
($-8.7\%$ to $-25\%$ comparisons) but caps the recall ceiling at 0.9817
(vs.\ 0.9947 at 1.2), so 1.2 is kept as the robust value.

\paragraph{O3: Exposed tuning levers.}
\texttt{prefetch\_lookahead} (sweep $\{8,16,24\}$: default 8 is already
optimal on the evaluation host) and \texttt{beam\_width} (beam $>1$ is a net
loss at 1T: +8--25\% comparisons, no recall gain) are exposed in
\texttt{job.json} as documented levers rather than committed wins.

\paragraph{O4: Fat LTO.}
Enabling \texttt{lto=fat} for release builds adds cross-crate inlining on the
call-heavy provider\,$\to$\,distance-kernel path: +10--12\% search QPS on a
10K A/B, identical comparisons/recall; release build time rises to
$\sim$6 minutes.

\section{Evaluation}
\label{sec:eval}
\subsection{Setup}
All cross-system numbers come from a single session (2026-09-16) on one
Tencent Cloud CVM instance: 2 vCPU, 3.6\,GB RAM + 8\,GB swap, OpenCloudOS 9.4.
Dataset SIFT1M (1M $\times$ 128, L2); recall@10; 10k queries; 3 repetitions,
median reported. Single-thread (1T) is the reference protocol; NSG's search
binary parallelizes queries over OpenMP with 2 threads and is labelled 2T.
Raw records ship in \texttt{bench/results/*.jsonl}.

\subsection{Main result: QPS at recall\,$\ge$\,0.95}
Table~\ref{tab:main} ranks systems at the recall\,$\ge$\,0.95 operating
point.

\begin{table}[t]
\centering
\small
\begin{tabular}{@{}lrrrrr@{}}
\toprule
System & QPS & recall@10 & Op.\ point & Build & Peak mem \\
\midrule
hnswlib-master-default & \textbf{6077} & 0.9637 & ef=64 & 238.6s & 1.31\,GB \\
hnswlib-stock (0.8.0) & 6017 & 0.9636 & ef=64 & 242.0s & 1.31\,GB \\
vamana (stock, $\alpha{=}1.0$) & 5708 & 0.9667 & Ls=64 & 255.0s & 1.27\,GB \\
hnswlib-variant4 & 5632 & 0.9617 & ef=64 & 260.4s & 1.31\,GB \\
diskann-rsi ($\alpha{=}1.0$) & 5078 & 0.9669 & Ls=64 & \textbf{131.9s} & 1.28\,GB \\
nsg (2T) & 2643 & 0.9525 & Ls=32 & 451s$^\dagger$ & 1.40\,GB \\
\bottomrule
\end{tabular}
\caption{QPS at recall\,$\ge 0.95$, single-thread unless noted.
$^\dagger$NSG additionally requires a kNN graph pre-computation of several
hours (efanna, K=200).}
\label{tab:main}
\end{table}

\subsection{Recall--QPS ladders}
Table~\ref{tab:ladder} gives the full operating ladders (1T; NSG 2T).
DiskANN-RSI trails stock Vamana by 31\% at Ls=16 and 19\% at Ls=32, reaches
$-11\%$ at Ls=64, and is within noise at Ls$\ge$128; recalls match to
$\pm$0.001 at every rung, indicating identical graph quality --- the delta is
pure search-kernel speed. hnswlib leads the $\ge$0.95 band on this host;
NSG trails throughout despite a 2-thread advantage.

\begin{table}[t]
\centering
\small
\begin{tabular}{@{}lllll@{}}
\toprule
System & \multicolumn{4}{l}{rung: recall/QPS (low $\to$ high)} \\
\midrule
vamana $\alpha{=}1.0$ & Ls16 .826/15382 & Ls32 .914/9705 & Ls64 .967/5708 & Ls128 .990/3178 \\
 & Ls256 .998/1756 & & & \\
vamana $\alpha{=}1.2$ & Ls16 .881/12513 & Ls32 .949/7853 & Ls64 .984/4534 & Ls128 .996/2658 \\
 & Ls256 .999/1448 & & & \\
diskann-rsi $\alpha{=}1.0$ & Ls16 .824/10587 & Ls32 .914/7897 & Ls64 .967/5078 & Ls128 .990/3090 \\
 & Ls256 .998/1719 & & & \\
diskann-rsi $\alpha{=}1.2$ & Ls16 .881/9398 & Ls32 .949/6650 & Ls64 .983/4164 & Ls128 .996/2502 \\
 & Ls256 .999/1435 & & & \\
hnswlib-stock & ef16 .801/16738 & ef32 .903/10281 & ef64 .964/6017 & ef128 .989/3279 \\
 & ef256 .997/1764 & & & \\
hnswlib-master-default & ef16 .801/16368 & ef32 .904/10448 & ef64 .964/6077 & ef128 .989/3258 \\
 & ef256 .997/1792 & & & \\
hnswlib-variant4 & ef16 .817/15621 & ef32 .907/9853 & ef64 .962/5632 & ef128 .986/3174 \\
 & ef256 .996/1728 & & & \\
nsg (2T) & Ls16 .877/2929 & Ls32 .952/2632 & Ls64 .986/2099 & Ls128 .996/1547 \\
 & Ls256 .999/1030 & & & \\
\bottomrule
\end{tabular}
\caption{Recall@10/QPS ladders on SIFT1M (1T; NSG 2T, medians of 3 reps).}
\label{tab:ladder}
\end{table}

\subsection{Thread scaling and build time}
At Ls=64, moving to two threads scales stock Vamana 1.75$\times$
(5708\,$\to$\,9993 QPS) and DiskANN-RSI 1.87$\times$ (5078\,$\to$\,9472);
at Ls=16 the ratios are 1.81$\times$ and 1.84$\times$. DiskANN-RSI builds
its index 1.9$\times$ faster than stock Vamana (131.9s vs.\ 255.0s) and
$1.8\times$ faster than hnswlib, at equal graph quality (matched recalls).

\subsection{Divergence from the optimizer's harness}
\label{sec:analysis}
The RSI harness reported +55\%; the unified benchmark shows $-9\%$ to $-31\%$
at low search depths and parity at high recall. Three factors explain the
gap. (i) \emph{Different protocols}: the RSI harness used its own data
layout, baseline pairing, and early-stop configuration; its baseline
(pre-optimization source) is not bit-comparable to the stock C++ Vamana
baseline used here. (ii) \emph{Hardware regime}: on this 2-vCPU host the 1M
working set is largely L3-resident, so the cache-residency win of 16-bit
stamps has little room to act, while the diversity-phase code added to the
kernel imposes constant per-expansion overhead that dominates at small L.
(iii) \emph{LTO asymmetry}: the RSI binary carries fat LTO (+10--12\%), i.e.,
the reported deficit is measured \emph{despite} an exclusive build-level
advantage. The honest summary is that O1--O4 are real but regime-dependent;
the durable, regime-independent wins are build time ($1.9\times$) and 2T
scaling.

\section{Negative Results}
\label{sec:negative}
The RSI loop also produced measured rejections, which we record to save
others the compute: \textbf{SQ8+rerank} tracks full precision within
$\sim$0.001 recall but yields no QPS gain here (5238 vs.\ 4641 at the metric
point) --- quantization pays when memory bandwidth is the bottleneck, which
an L3-resident working set is not; \textbf{candidate-queue merge insertion}
was slower and reverted; \textbf{random-projection data reordering} had no
measurable effect at 1M on this host; \textbf{partial-distance early exit}
aborts 86.9\% of distance computations but after $\sim$2.9 of 4 chunks, a
net loss for 128-dim vectors with a 4-chunk kernel; \textbf{beam width $>$1}
inflates comparisons 8--25\% with no recall gain at 1T;
\textbf{early-stop ratio 1.15} caps the recall ceiling at 0.9817.
Separately, an hnswlib pruning-heuristic variant (variant4) did not beat its
same-source default at the $\ge$0.95 band (5632 vs.\ 6077 QPS).

\section{Related Work}
\label{sec:related}
Beyond the graph-ANN literature of Section~\ref{sec:background}, this work
sits in automated research tooling: The AI Scientist
\cite{lu2024aiscientist} and v2 \cite{yamada2025aiscientistv2} automate idea
generation, experimentation, and writeup; Agent Laboratory
\cite{schmidgall2025agentlab} structures literature review, code solving,
and report writing as collaborating agents. Our RSI loop is of the same
family; this paper documents, with controlled re-measurement, the failure
mode of accepting optimizer-harness metrics at face value --- complementing
benchmarking methodology work such as ANN-Benchmarks
\cite{aumuller2020annbenchmarks}.

\section{Conclusion}
A Rust-native Vamana index with an automatically optimized query path wins
decisively inside its optimizer's harness (+55\%) and roughly ties stock
Vamana under unified same-machine conditions, with regime-independent wins
in build time ($1.9\times$) and two-thread scaling (1.87$\times$). The gap
between the two evaluations is itself a result: automated optimizers need
protocol-unified, cross-system validation gates before their improvements
are claimed as such. We release the code, the raw measurements, and the
runbook to make that gate cheap to apply.

\section*{Reproducibility}
Code, \texttt{job.json} configurations, raw JSONL measurements
(\texttt{bench/results/}), and the environment runbook are included in the
repository. The full optimization diff relative to the upstream port is in
\texttt{docs/final\_patch.diff}.

\begin{thebibliography}{9}

\bibitem{subramanya2019diskann}
S.~J. Subramanya, F.~Devvrit, R.~Kadekodi, R.~Krishnaswamy, H.~V. Simhadri.
\newblock DiskANN: Fast Accurate Billion-point Nearest Neighbor Search on a
Single Node.
\newblock \emph{NeurIPS}, Dec.\ 2019.

\bibitem{gollapudi2025sort}
S.~Gollapudi, R.~Krishnaswamy, K.~Shiragur, H.~Wardhan.
\newblock Sort Before You Prune: Improved Worst-Case Guarantees of the
DiskANN Family of Graphs.
\newblock \emph{ICML}, PMLR 267:19787--19808, Jul.\ 2025.

\bibitem{malkov2020hnsw}
Y.~A. Malkov, D.~A. Yashunin.
\newblock Efficient and Robust Approximate Nearest Neighbor Search Using
Hierarchical Navigable Small World Graphs.
\newblock \emph{IEEE TPAMI} 42(4):824--836, Apr.\ 2020. arXiv:1603.09320,
Mar.\ 2016.

\bibitem{fu2019nsg}
C.~Fu, C.~Xiang, C.~Wang, D.~Cai.
\newblock Fast Approximate Nearest Neighbor Search with the Navigating
Spreading-out Graph.
\newblock \emph{PVLDB} 12(5):461--474, Jan.\ 2019. arXiv:1707.00143, Jul.\ 2017.

\bibitem{indyk2023worstcase}
P.~Indyk, H.~Xu.
\newblock Worst-case Performance of Popular Approximate Nearest Neighbor
Search Implementations: Guarantees and Limitations.
\newblock \emph{NeurIPS Datasets and Benchmarks}, 2023. arXiv:2310.19126.

\bibitem{aumuller2020annbenchmarks}
M.~Aum\"uller, E.~Bernhardsson, A.~Faithfull.
\newblock ANN-Benchmarks: A Benchmarking Tool for Approximate Nearest
Neighbor Algorithms.
\newblock \emph{Information Systems} 87, 2020.

\bibitem{lu2024aiscientist}
C.~Lu, R.~T.~Q.~Chen, I.~Lillich, et al.
\newblock The AI Scientist: Towards Fully Automated Open-Ended Scientific
Discovery.
\newblock arXiv:2408.06292, Aug.\ 2024.

\bibitem{yamada2025aiscientistv2}
Y.~Yamada, R.~T.~Q.~Chen, et al.
\newblock The AI Scientist-v2: Workshop-Level Automated Scientific Discovery
via Agentic Tree Search.
\newblock arXiv:2504.08066, Apr.\ 2025.

\bibitem{schmidgall2025agentlab}
S.~Schmidgall, Y.~Su, Z.~Wang, et al.
\newblock Agent Laboratory: Using LLM Agents as Research Assistants.
\newblock arXiv:2501.04227, Jan.\ 2025.

\end{thebibliography}

\end{document}
